\sectionkicker{Source-backed technical notes}
\subsection*{生成メディアが独立した製品・モデルスタックへ: Technical Notes}
本欄は記事本文で圧縮した一次資料上の情報を、比較・再検証しやすい形で整理したものである。時系列、確認済みの主張、留意点、一次資料URLを掲載する。
\smallskip
\noindent{\small\textit{Technical Notesでは、一次資料から確認した公開・提供・研究上の事実を記載する。数値・能力評価や提供元・著者による比較表現は、独立再現済みの結果ではなく帰属claimとして扱う。}}\par
\medskip
\subsection*{Theme at a glance}
\addcontentsline{toc}{subsection}{Theme at a glance}
\begin{center}
\footnotesize
\begin{tabularx}{\linewidth}{@{}p{0.20\linewidth}p{0.10\linewidth}p{0.14\linewidth}X@{}}
\toprule
資料 & 位置づけ & 種別 & 時系列 \\
\midrule
DALL-E 3 2023 lifecycle & 主要資料 & モデル & — \\
Stable Diffusion XL / SDXL Turbo 2023 & 主要資料 & モデル更新 & — \\
Stable Video Diffusion 2023 & 主要資料 & モデル & — \\
Adobe Firefly / Generative Fill & 主要資料 & PRODUCT & — \\
\bottomrule
\end{tabularx}
\end{center}
\normalsize

\begin{technicalnote}{DALL-E 3 2023 lifecycle}{主要資料}
\begingroup
% reader-facing Technical Notes break policy
\widowpenalty=10000
\clubpenalty=10000
\displaywidowpenalty=10000
\begin{tabularx}{\linewidth}{@{}>{\bfseries}p{0.22\linewidth}X@{}}
組織 & OpenAI \\
種別 & モデル \\
時系列 & — \\
\end{tabularx}
\smallskip
{\bfseries 一次資料から整理したtechnical points}
\begin{itemize}[leftmargin=1.5em,itemsep=0.35em]
\item \textbf{一次情報で確認できる事実}: OpenAIの選定済み一次資料では、DALL·E 3 system card は2023年10月3日に公開され、model release に伴う safety / technical boundary を記録する一次資料として扱う。system-card publication と user-facing product availability は同一イベントではない。 2023年10月19日には DALL·E 3 が ChatGPT Plus と Enterprise で利用可能になったことが別途公表され、model / safety documentation から product availability への移行が確認できる。 OpenAI は wider release に向けた safety mitigation stack と provenance research の更新を説明しているが、その有効性評価は提供元の主張として扱う。
\end{itemize}
\begin{minipage}{\linewidth}
% reader-facing Technical Notes generic boundary/source tail group
\begin{samepage}
{\bfseries 一次資料}
\begingroup\sloppy
\Urlmuskip=0mu plus 2mu\relax
\begin{itemize}[leftmargin=1.5em,itemsep=0.25em]
\item {\scriptsize\url{https://openai.com/index/dall-e-3-is-now-available-in-chatgpt-plus-and-enterprise}}
\item {\scriptsize\url{https://openai.com/index/dall-e-3-system-card}}
\end{itemize}
\endgroup
\end{samepage}
\end{minipage}
\endgroup
\end{technicalnote}
\medskip

\begin{technicalnote}{Stable Diffusion XL / SDXL Turbo 2023}{主要資料}
\begingroup
% reader-facing Technical Notes break policy
\widowpenalty=10000
\clubpenalty=10000
\displaywidowpenalty=10000
\begin{tabularx}{\linewidth}{@{}>{\bfseries}p{0.22\linewidth}X@{}}
組織 & Stability AI \\
種別 & モデル更新 \\
時系列 & — \\
\end{tabularx}
\smallskip
{\bfseries 一次資料から整理したtechnical points}
\begin{itemize}[leftmargin=1.5em,itemsep=0.35em]
\item \textbf{一次情報で確認できる事実}: Stability AIの選定済み一次資料では、2023年4月の SDXL beta は API customers と DreamStudio users 向けに提供され、open-source release は当時まだ後続予定と明記されていたため、beta service availability と weights release を区別する。 2023年7月の SDXL 1.0 は 3.5B parameter base model と 6.6B parameter ensemble pipeline を持つ two-stage latent-diffusion architecture として説明され、base が noisy latents を生成した後に refinement model が最終 denoising を担う。weights / source code、Stability API、AWS、DreamStudio など複数の access surface が同時に示された。 SDXL 1.0 weights は CreativeML OpenRAIL++-M License で公開された一方、2023年11月の SDXL Turbo weights / code は non-commercial research license での公開であり、両者の license / commercial-use boundary を同一視しない。 SDXL Turbo は Adversarial Diffusion Distillation (ADD) を用いて single-step image generation を狙う model として公表され、対応論文は adversarial loss と score-distillation を組み合わせた 1--4 step sampling を記述する。性能比較値は Stability AI / 著者側の評価として扱う。
\end{itemize}
\begin{minipage}{\linewidth}
% reader-facing Technical Notes generic boundary/source tail group
\begin{samepage}
{\bfseries 一次資料}
\begingroup\sloppy
\Urlmuskip=0mu plus 2mu\relax
\begin{itemize}[leftmargin=1.5em,itemsep=0.25em]
\item {\scriptsize\url{http://arxiv.org/abs/2307.01952v1}}
\item {\scriptsize\url{http://arxiv.org/abs/2311.17042v1}}
\item {\scriptsize\url{https://stability.ai/news-updates/stability-ai-sdxl-turbo}}
\item {\scriptsize\url{https://stability.ai/news-updates/stable-diffusion-sdxl-1-announcement}}
\item {\scriptsize\url{https://stability.ai/news-updates/stable-diffusion-xl-beta-available-for-api-customers-and-dreamstudio-users}}
\end{itemize}
\endgroup
\end{samepage}
\end{minipage}
\endgroup
\end{technicalnote}
\medskip

\begin{technicalnote}{Stable Video Diffusion 2023}{主要資料}
\begingroup
% reader-facing Technical Notes break policy
\widowpenalty=10000
\clubpenalty=10000
\displaywidowpenalty=10000
\begin{tabularx}{\linewidth}{@{}>{\bfseries}p{0.22\linewidth}X@{}}
組織 & Stability AI \\
種別 & モデル \\
時系列 & — \\
\end{tabularx}
\smallskip
{\bfseries 一次資料から整理したtechnical points}
\begin{itemize}[leftmargin=1.5em,itemsep=0.35em]
\item \textbf{一次情報で確認できる事実}: Stability AIの選定済み一次資料では、2023年11月21日の Stable Video Diffusion は Stable Diffusion image model を基礎とする generative-video foundation model の research preview として公開され、code は GitHub、local execution 用 weights は Hugging Face で提供された。 research-preview 時点の公開物は image-to-video の2 model で、14 frames と 25 frames を 3--30 fps の範囲で生成すると説明された一方、real-world / commercial applications 向けではないと明記されていた。 2023年12月20日には Stable Video Diffusion が Stability AI Developer Platform API に追加され、programmatic image-to-video access が別の availability event として成立した。API版は 25 generated frames と 24 FILM-interpolated frames から2秒の動画を生成する product specification を持つ。 研究論文、research-preview weights、Developer Platform API はそれぞれ publication / open research access / hosted service の異なる lifecycle state として扱う。
\end{itemize}
\begin{minipage}{\linewidth}
% reader-facing Technical Notes generic boundary/source tail group
\begin{samepage}
{\bfseries 一次資料}
\begingroup\sloppy
\Urlmuskip=0mu plus 2mu\relax
\begin{itemize}[leftmargin=1.5em,itemsep=0.25em]
\item {\scriptsize\url{http://arxiv.org/abs/2311.15127v1}}
\item {\scriptsize\url{https://stability.ai/news-updates/introducing-stable-video-diffusion-api}}
\item {\scriptsize\url{https://stability.ai/news-updates/stable-video-diffusion-open-ai-video-model}}
\end{itemize}
\endgroup
\end{samepage}
\end{minipage}
\endgroup
\end{technicalnote}
\medskip

\begin{technicalnote}{Adobe Firefly / Generative Fill}{主要資料}
\begingroup
% reader-facing Technical Notes break policy
\widowpenalty=10000
\clubpenalty=10000
\displaywidowpenalty=10000
\begin{tabularx}{\linewidth}{@{}>{\bfseries}p{0.22\linewidth}X@{}}
組織 & Adobe \\
種別 & 製品 \\
時系列 & — \\
\end{tabularx}
\smallskip
{\bfseries 一次資料から整理したtechnical points}
\begin{itemize}[leftmargin=1.5em,itemsep=0.35em]
\item \textbf{一次情報で確認できる事実}: Adobeの選定済み一次資料では、2023年3月21日の Firefly 初回 beta は creative generative-AI model family の第一弾として image generation と text effects に焦点を置き、最初の model は Adobe Stock、openly licensed content、copyright-expired public-domain content を training source とする設計として説明された。 初回発表では Creative Cloud 等への将来 integration と API 提供計画が示されたが、それらを3月時点で既に一般提供済みだったとは扱わない。 2023年5月23日の Photoshop Generative Fill は Firefly を Photoshop workflow に直接統合し、text prompt から image content を add / extend / remove しつつ generated layers で non-destructive editing を行う beta feature として公開された。 同日、Generative Fill は Photoshop desktop beta と Firefly beta web module で利用可能とされ、Content Credentials 対応も説明された。将来の general availability と beta availability を区別する。
\end{itemize}
\begin{minipage}{\linewidth}
% reader-facing Technical Notes generic boundary/source tail group
\begin{samepage}
{\bfseries 一次資料}
\begingroup\sloppy
\Urlmuskip=0mu plus 2mu\relax
\begin{itemize}[leftmargin=1.5em,itemsep=0.25em]
\item {\scriptsize\url{https://news.adobe.com/news/news-details/2023/adobe-unveils-firefly-a-family-of-new-creative-generative-ai}}
\item {\scriptsize\url{https://news.adobe.com/news/news-details/2023/adobe-unveils-future-of-creative-cloud-with-generative-ai-as-a-creative-co-pilot-in-photoshop}}
\end{itemize}
\endgroup
\end{samepage}
\end{minipage}
\endgroup
\end{technicalnote}
\medskip
